\documentclass[a4paper,11pt]{article}

\usepackage[utf8]{inputenc}

\usepackage{minted}
\usepackage{graphicx}
\usepackage{placeins}

\begin{document}


\title{
    \textbf{T9 report}
}
\author{Jonathan Karlsson}
\date{Fall 2023}

\maketitle

\section*{Introduction}
The T9 assignment was to implement the T9 system, where a word is suggested depending on the keys that is inputted. The letters q and w were excluded in the assignment.
\section*{Implementation}
The T9 was implemented by first creating a constructor which only creates the root for the T9 as a new Node. Then four different 
support methods were created these were code, uncode, keyToInd and charToKey. Some repeating parts of the methods is cut out using three : to not take up too much space in this report.
The first method code took a char as argument and then using a switch statement returned 0 to 26 or -1 if the char was not in the standard swedish alphabet a...ö. 

\begin{minted}{java}
private static int code(char w) {
    switch (w) {
        case 'a' :
            return 0;
        case 'b' :
            return 1;
        case 'c' :
            return 2;
        case 'd' :
            return 3;
        :
        :
        :
        case 'ä' :
            return 25;
        case 'ö' :
            return 26;
    }
    return -1;
}    
\end{minted}
The second method was the uncode which did the opposite to the code method, taking in an int and returning a char depending on the value of the int. Then 
if the value was invalid, which means that it did not correspond to any letter the 0 char was returned.
\begin{minted}{java}
private static char uncode (int value) {
    switch(value) {
        case 0 :
            return 'a';
        case 1 :
            return 'b';
        case 2 :
            return 'c';
        case 3 :
            return 'd';
        :
        :
        :
        case 25 : 
            return 'ä';
        case 26 :
            return 'ö';
    }
    return 0;
}
\end{minted}
The third method takes a key character and returns an index from 0 to 8 depending on what key was inserted, this is so that a sequence of keys can be used as a string later on. 
If the key was not 0 to 9 then -1 is returned.
\begin{minted}{java}
private static int keyToInd(char key) {
    switch(key) {
        case '1' :
            return 0;
        :
        :
        :
        case '9' :
            return 8;
    }
    return -1;
}
\end{minted}
The final method, charToKey is used for testing so that a word can be turned into sequence of keys, this is why it is the only one of the support methods that is public. 
This method takes in a char and then returns key chars depending on what group the original character was in, groups mean "a,b,c","d,e,f" and so on. If 
the character was not in a group the 0 char is returned.
\begin{minted}{java}
public static char charToKey (char character) {
    switch(character) {
        case 'a':
        case 'b':
        case 'c':
            return '1';
        :
        :
        :
        case 'å':
        case 'ä':
        case 'ö':
            return '9';
    }
    return 0;
}
\end{minted}
There were two more methods in the T9 add and decode. The add method simply calls root with add using the word.
The decode method takes in a sequence then creates a new ArrayList, then calls colloct on the root using the sequence an empty string and the ArrayList. Finally the ArrayList is returned.
\begin{minted}{java}
public void add(String word) {
    root.add(word);
}
public ArrayList<String> decode(String sequence) {
    ArrayList<String> arrList = new ArrayList<String>();
    root.collect(sequence, "",arrList);
    return arrList;
}
\end{minted}
The T9 also contains a private class Node. This Node contains an array of Nodes called next, and a boolean called valid.
The node also contains three different methods the constructor, an add method and a collect method. 
The constructor simply creates the next array as an array of Nodes with size 27. It also sets valid to false.
\newline
The add method works by


\end{document}
